%\begin{slikaDesno}
\PID
Преносна карактеристика мерног система за мерење температуре описана је диференцијалном 
једначином првог реда у облику
$v + \uptau \dfrac{\de v}{\de t} = a(T - T_0)$, где је $v = v(t)$ напон на излазним прикључцима сензора, 
а $T = T(t)$ је температура коју он мери. Познате су вредности константи 
$a = 10\unit{mV/^\circ}$, $\uptau = 10\unit{s}$, и $T_0 = 25\unit{^\circ C}$.
\begin{enumerate}[label=(\alph*)]
\item Сензор се налази на амбијентној температури $T_0$, а у тренутку $t = 0$ се нагло урања 
у течност температуре $T = 80\unit{^\circ C}$.  
Одредити временски облик напон на излазним крајевима након потапања у течност $v(t > 0)$. 
\item Одредити време успона, $t_{\rm r}$ (\textit{rise time}), за које се одзив сензора промени од $10\%$ до 
$90\%$ укупне промене током процеса из претходне тачке.  
\item Израчунати пропусни опсег сензора, $f_{\rm kr}$, ($3\unit{dB}$ пропсуни опсег), на основу задатих параметара. 
\end{enumerate}

